\documentclass[aspectratio=169, 11pt]{beamer}
% ============================================================================
% Preambulo compartido para presentaciones Beamer
% Estadistica 2026-II - Pontificia Universidad Javeriana
% Incluir con: % ============================================================================
% Preambulo compartido para presentaciones Beamer
% Estadistica 2026-II - Pontificia Universidad Javeriana
% Incluir con: % ============================================================================
% Preambulo compartido para presentaciones Beamer
% Estadistica 2026-II - Pontificia Universidad Javeriana
% Incluir con: \input{../preambulo_beamer.tex} despues de \documentclass
% ============================================================================

\usepackage[T1]{fontenc}
\usepackage[utf8]{inputenc}
\usepackage[spanish]{babel}
\usepackage{amsmath, amssymb, amsfonts}
\usepackage{booktabs}
\usepackage{graphicx}
\usepackage{tikz}
\usetikzlibrary{shapes.geometric, positioning, arrows.meta, calc}
\usepackage{pgfplots}
\pgfplotsset{compat=1.18}
\usepackage{listings}
\usepackage{xcolor}
\usepackage{hyperref}
\usepackage{multicol}

% --- Tema Beamer (minimalista) ---
\usetheme{default}
\usecolortheme{default}
\usefonttheme{default}
\setbeamertemplate{navigation symbols}{}
\setbeamertemplate{caption}[numbered]
\setbeamertemplate{footline}{%
  \hspace{1em}{\tiny\url{https://github.com/polanco-jaime/Curso-Estadistica-2026-3}}\hfill\usebeamerfont{footline}\insertframenumber/\inserttotalframenumber\hspace{1em}\vspace{0.5em}%
}
\setbeamertemplate{itemize items}[circle]
\setbeamertemplate{enumerate items}[default]

% --- Colores institucionales (sutiles) ---
\definecolor{javeriana}{RGB}{0, 62, 105}
\definecolor{javerianalight}{RGB}{0, 105, 170}
\definecolor{codigobg}{RGB}{245, 245, 245}
\definecolor{codigofg}{RGB}{40, 40, 40}
\definecolor{comentario}{RGB}{100, 100, 100}
\definecolor{keyword}{RGB}{0, 100, 150}
\definecolor{string}{RGB}{180, 60, 30}

\setbeamercolor{title}{fg=javeriana}
\setbeamercolor{frametitle}{fg=javeriana}
\setbeamercolor{structure}{fg=javeriana}
\setbeamercolor{block title}{fg=white, bg=javeriana!75}
\setbeamercolor{block body}{bg=javeriana!5}
\setbeamercolor{block title alerted}{fg=white, bg=red!70!black}
\setbeamercolor{block body alerted}{bg=red!5}
\setbeamercolor{block title example}{fg=white, bg=green!40!black}
\setbeamercolor{block body example}{bg=green!5}
\setbeamercolor{item}{fg=javeriana}

\setbeamerfont{title}{series=\bfseries, size=\Large}
\setbeamerfont{frametitle}{series=\bfseries}

% --- Configuracion de listings para R ---
\lstset{
  language=R,
  basicstyle=\ttfamily\footnotesize,
  backgroundcolor=\color{codigobg},
  commentstyle=\color{comentario}\itshape,
  keywordstyle=\color{keyword}\bfseries,
  stringstyle=\color{string},
  numbers=none,
  frame=single,
  rulecolor=\color{javeriana!30},
  breaklines=true,
  showstringspaces=false,
  columns=flexible,
  morekeywords={library, ggplot, aes, geom_histogram, geom_boxplot,
    geom_point, geom_smooth, geom_density, geom_bar, geom_line,
    geom_vline, geom_hline, geom_errorbar, labs, theme_minimal,
    facet_wrap, scale_fill_manual, scale_color_manual,
    summarise, group_by, filter, mutate, select, arrange,
    read.csv, head, str, summary, mean, median, sd, var, cor,
    lm, t.test, chisq.test, wilcox.test, kruskal.test, aov,
    pnorm, qnorm, dnorm, rnorm, qt, pt, qchisq, pchisq, qf, pf,
    confint, coeftest, vcovHC, bptest, vif, cooks.distance,
    TukeyHSD, table, prop.table}
}

% --- Comandos personalizados ---
\newcommand{\R}{\texttt{R}}
\newcommand{\code}[1]{\texttt{\footnotesize #1}}
\newcommand{\variable}[1]{\texttt{#1}}
\newcommand{\paquete}[1]{\texttt{#1}}

% --- Informacion del curso ---
\institute{Pontificia Universidad Javeriana\\
  Facultad de Ciencias Econ\'omicas y Administrativas}
\date{2026-II}
 despues de \documentclass
% ============================================================================

\usepackage[T1]{fontenc}
\usepackage[utf8]{inputenc}
\usepackage[spanish]{babel}
\usepackage{amsmath, amssymb, amsfonts}
\usepackage{booktabs}
\usepackage{graphicx}
\usepackage{tikz}
\usetikzlibrary{shapes.geometric, positioning, arrows.meta, calc}
\usepackage{pgfplots}
\pgfplotsset{compat=1.18}
\usepackage{listings}
\usepackage{xcolor}
\usepackage{hyperref}
\usepackage{multicol}

% --- Tema Beamer (minimalista) ---
\usetheme{default}
\usecolortheme{default}
\usefonttheme{default}
\setbeamertemplate{navigation symbols}{}
\setbeamertemplate{caption}[numbered]
\setbeamertemplate{footline}{%
  \hspace{1em}{\tiny\url{https://github.com/polanco-jaime/Curso-Estadistica-2026-3}}\hfill\usebeamerfont{footline}\insertframenumber/\inserttotalframenumber\hspace{1em}\vspace{0.5em}%
}
\setbeamertemplate{itemize items}[circle]
\setbeamertemplate{enumerate items}[default]

% --- Colores institucionales (sutiles) ---
\definecolor{javeriana}{RGB}{0, 62, 105}
\definecolor{javerianalight}{RGB}{0, 105, 170}
\definecolor{codigobg}{RGB}{245, 245, 245}
\definecolor{codigofg}{RGB}{40, 40, 40}
\definecolor{comentario}{RGB}{100, 100, 100}
\definecolor{keyword}{RGB}{0, 100, 150}
\definecolor{string}{RGB}{180, 60, 30}

\setbeamercolor{title}{fg=javeriana}
\setbeamercolor{frametitle}{fg=javeriana}
\setbeamercolor{structure}{fg=javeriana}
\setbeamercolor{block title}{fg=white, bg=javeriana!75}
\setbeamercolor{block body}{bg=javeriana!5}
\setbeamercolor{block title alerted}{fg=white, bg=red!70!black}
\setbeamercolor{block body alerted}{bg=red!5}
\setbeamercolor{block title example}{fg=white, bg=green!40!black}
\setbeamercolor{block body example}{bg=green!5}
\setbeamercolor{item}{fg=javeriana}

\setbeamerfont{title}{series=\bfseries, size=\Large}
\setbeamerfont{frametitle}{series=\bfseries}

% --- Configuracion de listings para R ---
\lstset{
  language=R,
  basicstyle=\ttfamily\footnotesize,
  backgroundcolor=\color{codigobg},
  commentstyle=\color{comentario}\itshape,
  keywordstyle=\color{keyword}\bfseries,
  stringstyle=\color{string},
  numbers=none,
  frame=single,
  rulecolor=\color{javeriana!30},
  breaklines=true,
  showstringspaces=false,
  columns=flexible,
  morekeywords={library, ggplot, aes, geom_histogram, geom_boxplot,
    geom_point, geom_smooth, geom_density, geom_bar, geom_line,
    geom_vline, geom_hline, geom_errorbar, labs, theme_minimal,
    facet_wrap, scale_fill_manual, scale_color_manual,
    summarise, group_by, filter, mutate, select, arrange,
    read.csv, head, str, summary, mean, median, sd, var, cor,
    lm, t.test, chisq.test, wilcox.test, kruskal.test, aov,
    pnorm, qnorm, dnorm, rnorm, qt, pt, qchisq, pchisq, qf, pf,
    confint, coeftest, vcovHC, bptest, vif, cooks.distance,
    TukeyHSD, table, prop.table}
}

% --- Comandos personalizados ---
\newcommand{\R}{\texttt{R}}
\newcommand{\code}[1]{\texttt{\footnotesize #1}}
\newcommand{\variable}[1]{\texttt{#1}}
\newcommand{\paquete}[1]{\texttt{#1}}

% --- Informacion del curso ---
\institute{Pontificia Universidad Javeriana\\
  Facultad de Ciencias Econ\'omicas y Administrativas}
\date{2026-II}
 despues de \documentclass
% ============================================================================

\usepackage[T1]{fontenc}
\usepackage[utf8]{inputenc}
\usepackage[spanish]{babel}
\usepackage{amsmath, amssymb, amsfonts}
\usepackage{booktabs}
\usepackage{graphicx}
\usepackage{tikz}
\usetikzlibrary{shapes.geometric, positioning, arrows.meta, calc}
\usepackage{pgfplots}
\pgfplotsset{compat=1.18}
\usepackage{listings}
\usepackage{xcolor}
\usepackage{hyperref}
\usepackage{multicol}

% --- Tema Beamer (minimalista) ---
\usetheme{default}
\usecolortheme{default}
\usefonttheme{default}
\setbeamertemplate{navigation symbols}{}
\setbeamertemplate{caption}[numbered]
\setbeamertemplate{footline}{%
  \hspace{1em}{\tiny\url{https://github.com/polanco-jaime/Curso-Estadistica-2026-3}}\hfill\usebeamerfont{footline}\insertframenumber/\inserttotalframenumber\hspace{1em}\vspace{0.5em}%
}
\setbeamertemplate{itemize items}[circle]
\setbeamertemplate{enumerate items}[default]

% --- Colores institucionales (sutiles) ---
\definecolor{javeriana}{RGB}{0, 62, 105}
\definecolor{javerianalight}{RGB}{0, 105, 170}
\definecolor{codigobg}{RGB}{245, 245, 245}
\definecolor{codigofg}{RGB}{40, 40, 40}
\definecolor{comentario}{RGB}{100, 100, 100}
\definecolor{keyword}{RGB}{0, 100, 150}
\definecolor{string}{RGB}{180, 60, 30}

\setbeamercolor{title}{fg=javeriana}
\setbeamercolor{frametitle}{fg=javeriana}
\setbeamercolor{structure}{fg=javeriana}
\setbeamercolor{block title}{fg=white, bg=javeriana!75}
\setbeamercolor{block body}{bg=javeriana!5}
\setbeamercolor{block title alerted}{fg=white, bg=red!70!black}
\setbeamercolor{block body alerted}{bg=red!5}
\setbeamercolor{block title example}{fg=white, bg=green!40!black}
\setbeamercolor{block body example}{bg=green!5}
\setbeamercolor{item}{fg=javeriana}

\setbeamerfont{title}{series=\bfseries, size=\Large}
\setbeamerfont{frametitle}{series=\bfseries}

% --- Configuracion de listings para R ---
\lstset{
  language=R,
  basicstyle=\ttfamily\footnotesize,
  backgroundcolor=\color{codigobg},
  commentstyle=\color{comentario}\itshape,
  keywordstyle=\color{keyword}\bfseries,
  stringstyle=\color{string},
  numbers=none,
  frame=single,
  rulecolor=\color{javeriana!30},
  breaklines=true,
  showstringspaces=false,
  columns=flexible,
  morekeywords={library, ggplot, aes, geom_histogram, geom_boxplot,
    geom_point, geom_smooth, geom_density, geom_bar, geom_line,
    geom_vline, geom_hline, geom_errorbar, labs, theme_minimal,
    facet_wrap, scale_fill_manual, scale_color_manual,
    summarise, group_by, filter, mutate, select, arrange,
    read.csv, head, str, summary, mean, median, sd, var, cor,
    lm, t.test, chisq.test, wilcox.test, kruskal.test, aov,
    pnorm, qnorm, dnorm, rnorm, qt, pt, qchisq, pchisq, qf, pf,
    confint, coeftest, vcovHC, bptest, vif, cooks.distance,
    TukeyHSD, table, prop.table}
}

% --- Comandos personalizados ---
\newcommand{\R}{\texttt{R}}
\newcommand{\code}[1]{\texttt{\footnotesize #1}}
\newcommand{\variable}[1]{\texttt{#1}}
\newcommand{\paquete}[1]{\texttt{#1}}

% --- Informacion del curso ---
\institute{Pontificia Universidad Javeriana\\
  Facultad de Ciencias Econ\'omicas y Administrativas}
\date{2026-II}

\title{Sesión 7: Pruebas para Dos Poblaciones, ANOVA y No Paramétricas}
\subtitle{Comparando grupos: de la prueba t a Kruskal-Wallis}
\author{Prof.\ Jaime Polanco Jim\'enez}
\date{Viernes 2 de octubre de 2026}
\begin{document}

\begin{frame}
\titlepage
\end{frame}

\begin{frame}{Agenda}
\tableofcontents
\end{frame}

% ============================================================================
\section{Prueba t Pareada (Muestras Dependientes)}
% ============================================================================

\begin{frame}{Muestras Independientes vs Dependientes}
\textbf{Hasta ahora}: Prueba $t$ para dos muestras \alert{independientes}
\begin{itemize}
  \item Grupos completamente distintos (ej. públicas vs privadas)
  \item Sin relación entre observaciones de un grupo y otro
\end{itemize}

\vspace{0.3cm}
\textbf{Ahora}: Muestras \alert{dependientes (pareadas)}
\begin{itemize}
  \item Mismos sujetos medidos dos veces (antes/después)
  \item Sujetos emparejados naturalmente (hermanos, pares matched)
  \item Cada observación en el grupo 1 tiene una contraparte específica en el grupo 2
\end{itemize}

\vspace{0.3cm}
\textbf{Ejemplos}:
\begin{itemize}
  \item Puntaje Saber 11 vs Saber Pro del \alert{mismo} estudiante
  \item Rendimiento antes vs después de una intervención
  \item Calificación de dos evaluadores sobre el \alert{mismo} trabajo
\end{itemize}
\end{frame}

\begin{frame}{¿Cuándo Usar Prueba t Pareada?}
\textbf{Criterios}:
\begin{enumerate}
  \item Dos mediciones del \alert{mismo} sujeto/unidad
  \item O pares naturales (matched pairs)
  \item Las diferencias individuales importan
\end{enumerate}

\vspace{0.3cm}
\textbf{Ventaja}:
\begin{itemize}
  \item Controla la variabilidad entre sujetos
  \item Mayor \alert{potencia} que la prueba de dos muestras independientes
  \item Reduce el error tipo II
\end{itemize}

\vspace{0.3cm}
\textbf{Ejemplo ICFES}:
\begin{itemize}
  \item Comparar inglés Saber 11 vs Saber Pro del mismo estudiante
  \item Problema: los datos ICFES actuales NO tienen esta vinculación (cada fila = un estudiante en un momento)
  \item Solución: usar datos simulados o buscar tabla de cruce
\end{itemize}
\end{frame}

\begin{frame}{Hipótesis y Estadístico de Prueba}
\textbf{Enfoque}: Trabajar con las \alert{diferencias} $d_i = x_{1i} - x_{2i}$

\vspace{0.3cm}
\textbf{Hipótesis}:
\[
H_0: \mu_d = 0 \quad \text{vs} \quad H_1: \mu_d \neq 0
\]
donde $\mu_d$ = media poblacional de las diferencias.

\vspace{0.3cm}
\begin{block}{Estadístico de prueba}
\[
t = \frac{\bar{d}}{s_d / \sqrt{n}}
\]
donde:
\begin{itemize}
  \item $\bar{d} = \frac{1}{n} \sum_{i=1}^n d_i$ = media de las diferencias
  \item $s_d$ = desviación estándar de las diferencias
  \item $n$ = número de pares
  \item $t \sim t(n-1)$ si $H_0$ es verdadera
\end{itemize}
\end{block}
\end{frame}

\begin{frame}[fragile]{Ejemplo: Saber 11 vs Saber Pro (Simulado) - Parte 1}
\begin{lstlisting}
# Simular datos pareados (en la practica, vienen de una tabla de cruce)
set.seed(123)
n <- 80
saber11 <- rnorm(n, mean = 155, sd = 18)
saber_pro <- saber11 + rnorm(n, mean = 8, sd = 5)
# Saber Pro suele ser ligeramente mayor

# Calcular diferencias
diferencias <- saber_pro - saber11

# Descriptivos
cat("Media Saber 11:", round(mean(saber11), 2))
cat("\nMedia Saber Pro:", round(mean(saber_pro), 2))
cat("\nMedia de diferencias:", round(mean(diferencias), 2))
cat("\nDE de diferencias:", round(sd(diferencias), 2))
\end{lstlisting}
\end{frame}

\begin{frame}[fragile]{Ejemplo: Saber 11 vs Saber Pro (Simulado) - Parte 2}
\begin{lstlisting}
# Visualizar diferencias
hist(diferencias, breaks = 15, col = "steelblue",
     main = "Distribucion de las Diferencias",
     xlab = "Saber Pro - Saber 11")
abline(v = 0, col = "red", lwd = 2, lty = 2)
\end{lstlisting}
\end{frame}

\begin{frame}[fragile]{Realizar la Prueba t Pareada}
\begin{lstlisting}
# Metodo 1: t.test con paired = TRUE
resultado <- t.test(saber_pro, saber11, paired = TRUE)

print(resultado)

cat("\nEstadistico t:", resultado$statistic)
cat("\np-valor:", resultado$p.value)
cat("\nIC para mu_d:", resultado$conf.int)

# Metodo 2: t.test sobre las diferencias directamente
resultado2 <- t.test(diferencias, mu = 0)
# Es equivalente al metodo 1

# Decision
if (resultado$p.value < 0.05) {
  cat("\nConcluimos que hay diferencia significativa")
  cat("\nSaber Pro difiere de Saber 11")
} else {
  cat("\nNo hay evidencia de diferencia")
}
\end{lstlisting}
\end{frame}

\begin{frame}[fragile]{Visualización de Datos Pareados}
\begin{lstlisting}
library(ggplot2)

# Crear data frame
datos_pareados <- data.frame(
  id = rep(1:n, 2),
  examen = rep(c("Saber 11", "Saber Pro"), each = n),
  puntaje = c(saber11, saber_pro)
)

# Grafico de lineas conectando pares
ggplot(datos_pareados, aes(x = examen, y = puntaje, group = id)) +
  geom_line(alpha = 0.3, color = "gray50") +
  geom_point(alpha = 0.5) +
  stat_summary(aes(group = 1), fun = mean, geom = "line",
               color = "red", size = 1.5) +
  stat_summary(fun = mean, geom = "point",
               color = "red", size = 4) +
  labs(title = "Comparacion Pareada: Saber 11 vs Saber Pro",
       y = "Puntaje de Ingles") +
  theme_minimal()
\end{lstlisting}
\end{frame}

% ============================================================================
\section{ANOVA de Una Vía}
% ============================================================================

\begin{frame}{Motivación: Comparar Más de Dos Grupos}
\textbf{Problema}: Queremos comparar inglés entre 5 regiones de Colombia.

\vspace{0.3cm}
\textbf{Opción incorrecta}: Hacer múltiples pruebas $t$
\begin{itemize}
  \item Comparar Caribe vs Andina, Caribe vs Pacífico, ..., etc.
  \item Con $k = 5$ grupos, tendríamos $\binom{5}{2} = 10$ comparaciones
  \item Problema: \alert{inflación del error tipo I}
\end{itemize}

\vspace{0.3cm}
\begin{alertblock}{Inflación del error tipo I}
Si hacemos $m$ pruebas independientes al nivel $\alpha$:
\[
P(\text{al menos un error tipo I}) = 1 - (1-\alpha)^m
\]
Ejemplo: 10 pruebas con $\alpha = 0.05$:
\[
P(\text{al menos un error}) = 1 - 0.95^{10} = 0.40 \quad (40\%!)
\]
\end{alertblock}

\vspace{0.3cm}
\textbf{Solución}: ANOVA (Analysis of Variance)
\end{frame}

\begin{frame}{¿Qué es ANOVA?}
\begin{definition}[ANOVA de una vía]
Una prueba que compara las medias de \alert{tres o más} grupos simultáneamente, controlando el error tipo I global.
\end{definition}

\vspace{0.3cm}
\textbf{Hipótesis}:
\[
H_0: \mu_1 = \mu_2 = \cdots = \mu_k \quad \text{vs} \quad H_1: \text{al menos un par difiere}
\]

\vspace{0.3cm}
\textbf{Importante}: $H_1$ NO dice \emph{cuáles} grupos difieren, solo que \alert{al menos uno} es distinto.

\vspace{0.3cm}
\textbf{Nombre}: ``Analysis of Variance'' porque comparamos variabilidad \alert{entre} grupos vs variabilidad \alert{dentro} de grupos.
\end{frame}

\begin{frame}{Lógica de ANOVA}
\textbf{Idea central}:
\begin{itemize}
  \item Si $H_0$ es verdadera (todas las medias son iguales), la variabilidad entre grupos debería ser similar a la variabilidad dentro de grupos (solo ruido aleatorio)
  \item Si $H_0$ es falsa (medias diferentes), la variabilidad entre grupos será \alert{mayor} que la variabilidad dentro de grupos
\end{itemize}

\vspace{0.3cm}
\begin{block}{Razón F}
\[
F = \frac{\text{Variabilidad entre grupos}}{\text{Variabilidad dentro de grupos}} = \frac{\text{MSB}}{\text{MSW}}
\]
\begin{itemize}
  \item Si $F$ es grande $\Rightarrow$ evidencia contra $H_0$
  \item Si $H_0$ es verdadera, $F \sim F(k-1, N-k)$ (distribución $F$)
\end{itemize}
\end{block}
\end{frame}

\begin{frame}{Partición de la Variabilidad Total}
\textbf{Suma de cuadrados total (SST)}:
\[
\text{SST} = \sum_{i=1}^k \sum_{j=1}^{n_i} (x_{ij} - \bar{x})^2
\]
Mide la variabilidad total de todos los datos.

\vspace{0.3cm}
\textbf{Descomposición}:
\[
\text{SST} = \text{SSB} + \text{SSW}
\]
\begin{itemize}
  \item \textbf{SSB (between)}: $\sum_{i=1}^k n_i (\bar{x}_i - \bar{x})^2$ = variabilidad entre grupos
  \item \textbf{SSW (within)}: $\sum_{i=1}^k \sum_{j=1}^{n_i} (x_{ij} - \bar{x}_i)^2$ = variabilidad dentro de grupos
\end{itemize}

\vspace{0.3cm}
\textbf{Cuadrados medios}:
\[
\text{MSB} = \frac{\text{SSB}}{k-1}, \quad \text{MSW} = \frac{\text{SSW}}{N-k}
\]
\end{frame}

\begin{frame}{Tabla ANOVA}
\begin{table}
\centering
\begin{tabular}{lcccc}
\toprule
\textbf{Fuente} & \textbf{SS} & \textbf{df} & \textbf{MS} & \textbf{F} \\
\midrule
Entre grupos & SSB & $k-1$ & $\frac{\text{SSB}}{k-1}$ & $\frac{\text{MSB}}{\text{MSW}}$ \\[0.2cm]
Dentro grupos & SSW & $N-k$ & $\frac{\text{SSW}}{N-k}$ & \\[0.2cm]
Total & SST & $N-1$ & & \\
\bottomrule
\end{tabular}
\end{table}

\vspace{0.3cm}
\textbf{Notación}:
\begin{itemize}
  \item $k$ = número de grupos
  \item $N = \sum_{i=1}^k n_i$ = tamaño total de la muestra
  \item $n_i$ = tamaño del grupo $i$
\end{itemize}

\vspace{0.3cm}
\textbf{p-valor}: $P(F_{k-1, N-k} > F_{\text{obs}})$
\end{frame}

\begin{frame}[fragile]{Ejemplo: Inglés por Macro-Región - Parte 1}
\small
\begin{lstlisting}
# Crear variable de macro-region (simplificacion)
icfes_ni$MACRO_REGION <- ifelse(
  icfes_ni$ESTU_DEPTO_RESIDE %in%
    c("Atlantico", "Bolivar", "Cesar", "Cordoba", "La Guajira",
      "Magdalena", "Sucre", "San Andres"),
  "Caribe",
  ifelse(
    icfes_ni$ESTU_DEPTO_RESIDE %in%
      c("Antioquia", "Bogota", "Boyaca", "Caldas", "Cundinamarca",
        "Huila", "Norte Santander", "Quindio", "Risaralda",
        "Santander", "Tolima"),
    "Andina",
\end{lstlisting}
\end{frame}

\begin{frame}[fragile]{Ejemplo: Inglés por Macro-Región - Parte 2}
\small
\begin{lstlisting}
# Continuacion: macro-region
    ifelse(
      icfes_ni$ESTU_DEPTO_RESIDE %in%
        c("Choco", "Cauca", "Narino", "Valle"),
      "Pacifica",
      ifelse(
        icfes_ni$ESTU_DEPTO_RESIDE %in%
          c("Amazonas", "Caqueta", "Guainia", "Guaviare",
            "Meta", "Putumayo", "Vaupes", "Vichada"),
        "Orinoquia-Amazonia", "Otras"
      )
    )
  )
)
\end{lstlisting}
\end{frame}

\begin{frame}[fragile]{Realizar ANOVA en R}
\begin{lstlisting}
# Estadisticas descriptivas por region
library(dplyr)
desc_regiones <- icfes_ni %>%
  group_by(MACRO_REGION) %>%
  summarise(
    n = n(),
    media = mean(MOD_INGLES_PUNT, na.rm = TRUE),
    sd = sd(MOD_INGLES_PUNT, na.rm = TRUE)
  )
print(desc_regiones)

# ANOVA
modelo_anova <- aov(MOD_INGLES_PUNT ~ MACRO_REGION, data = icfes_ni)

# Ver tabla ANOVA
summary(modelo_anova)
\end{lstlisting}
\end{frame}

\begin{frame}[fragile]{Interpretar Output de ANOVA}
\begin{lstlisting}[basicstyle=\ttfamily\footnotesize]
                  Df Sum Sq Mean Sq F value   Pr(>F)
MACRO_REGION       4  12458  3114.5   9.876 1.24e-07 ***
Residuals        495 156123   315.4
---
Signif. codes:  0 '***' 0.001 '**' 0.01 '*' 0.05 '.' 0.1 ' ' 1
\end{lstlisting}

\vspace{0.3cm}
\textbf{Interpretación}:
\begin{itemize}
  \item $F = 9.876$ con gl = $(4, 495)$
  \item p-valor $< 0.001$ $\Rightarrow$ \alert{Rechazamos $H_0$}
  \item Conclusión: Hay evidencia de que al menos una macro-región tiene media diferente de las demás
  \item \textbf{Pregunta pendiente}: ¿Cuáles regiones difieren? (Tukey HSD)
\end{itemize}
\end{frame}

\begin{frame}{Supuestos de ANOVA}
\textbf{Para que ANOVA sea válida, se requiere}:

\vspace{0.3cm}
\begin{enumerate}
  \item \textbf{Normalidad}: Los residuos deben seguir una distribución normal
  \begin{itemize}
    \item Verificar con QQ-plot o Shapiro-Wilk
    \item Robusta si $n$ es grande (TLC)
  \end{itemize}

  \item \textbf{Homocedasticidad}: Las varianzas de los grupos deben ser iguales
  \begin{itemize}
    \item Verificar con prueba de Levene o gráfico de residuos
    \item Regla empírica: $\max(s_i) / \min(s_i) < 2$
  \end{itemize}

  \item \textbf{Independencia}: Las observaciones deben ser independientes
  \begin{itemize}
    \item Garantizada por diseño de muestreo
  \end{itemize}
\end{enumerate}

\vspace{0.3cm}
\textbf{Si los supuestos no se cumplen}: Considerar transformación de datos o prueba no paramétrica (Kruskal-Wallis).
\end{frame}

\begin{frame}[fragile]{Verificar Supuestos - Parte 1}
\begin{lstlisting}
# 1. Grafico QQ de residuos
qqnorm(residuals(modelo_anova))
qqline(residuals(modelo_anova), col = "red")

# 2. Prueba de Shapiro-Wilk (si n no es muy grande)
shapiro.test(residuals(modelo_anova)[1:5000])
# Nota: Shapiro-Wilk tiene limite de n

# 3. Prueba de Levene para homocedasticidad
library(car)
leveneTest(MOD_INGLES_PUNT ~ MACRO_REGION, data = icfes_ni)

# 4. Grafico de residuos vs valores ajustados
plot(modelo_anova, which = 1)
# Debe verse sin patron sistematico
\end{lstlisting}
\end{frame}

\begin{frame}[fragile]{Verificar Supuestos - Parte 2}
\begin{lstlisting}
# 5. Boxplot por grupo (visualizar homogeneidad)
boxplot(MOD_INGLES_PUNT ~ MACRO_REGION, data = icfes_ni,
        col = "lightblue", ylab = "Puntaje Ingles",
        main = "Distribucion por Macro-Region")
\end{lstlisting}
\end{frame}

% ============================================================================
\section{Comparaciones Múltiples: Tukey HSD}
% ============================================================================

\begin{frame}{El Problema de las Comparaciones Múltiples}
\textbf{Situación}: ANOVA nos dice que HAY diferencias, pero no DÓNDE.

\vspace{0.3cm}
\textbf{¿Podemos hacer pruebas $t$ entre todos los pares?}
\begin{itemize}
  \item Con $k = 5$ grupos: $\binom{5}{2} = 10$ comparaciones
  \item Inflación del error tipo I (como vimos antes)
\end{itemize}

\vspace{0.3cm}
\textbf{Solución}: Métodos de comparaciones múltiples que controlan la \alert{tasa de error familiar} (FWER: Family-Wise Error Rate).

\vspace{0.3cm}
\textbf{Métodos comunes}:
\begin{itemize}
  \item \textbf{Tukey HSD} (Honestly Significant Difference): compara todos los pares
  \item \textbf{Bonferroni}: ajusta $\alpha$ dividiéndolo por el número de comparaciones
  \item \textbf{Scheffé}: más conservador, útil para comparaciones complejas
  \item \textbf{Dunnett}: compara todos vs un control
\end{itemize}

\vspace{0.3cm}
\textbf{Recomendación}: Tukey HSD es el más usado para comparaciones post-hoc.
\end{frame}

\begin{frame}{Tukey HSD: Honestly Significant Difference}
\textbf{Idea}: Ajustar el valor crítico para tener en cuenta las múltiples comparaciones.

\vspace{0.3cm}
\begin{block}{Intervalo de confianza simultáneo al $(1-\alpha)$}
Para cada par $(i, j)$:
\[
(\bar{x}_i - \bar{x}_j) \pm q_{\alpha, k, N-k} \cdot \sqrt{\frac{\text{MSW}}{2} \left(\frac{1}{n_i} + \frac{1}{n_j}\right)}
\]
donde $q_{\alpha, k, N-k}$ es el valor crítico de la distribución \alert{studentized range}.
\end{block}

\vspace{0.3cm}
\textbf{Interpretación}:
\begin{itemize}
  \item Si el IC NO incluye 0, los grupos $i$ y $j$ difieren significativamente
  \item Controla FWER al nivel $\alpha$ para \alert{todas} las comparaciones simultáneamente
\end{itemize}
\end{frame}

\begin{frame}[fragile]{Realizar Tukey HSD en R}
\begin{lstlisting}
# Prueba de Tukey HSD
tukey_resultado <- TukeyHSD(modelo_anova, conf.level = 0.95)

print(tukey_resultado)

# Visualizar resultados
plot(tukey_resultado, las = 1, cex.axis = 0.7)
abline(v = 0, col = "red", lty = 2)
\end{lstlisting}

\vspace{0.3cm}
\textbf{Output}: Tabla con todas las comparaciones pareadas
\begin{itemize}
  \item \textbf{diff}: Diferencia de medias entre los grupos
  \item \textbf{lwr, upr}: Límites del IC ajustado
  \item \textbf{p adj}: p-valor ajustado
\end{itemize}
\end{frame}

\begin{frame}[fragile]{Interpretar Resultados de Tukey}
\begin{lstlisting}[basicstyle=\ttfamily\tiny]
  Tukey multiple comparisons of means
    95% family-wise confidence level

Fit: aov(formula = MOD_INGLES_PUNT ~ MACRO_REGION, data = icfes_ni)

$MACRO_REGION
                             diff       lwr       upr     p adj
Caribe-Andina            -8.35421 -14.2893  -2.41915  0.001823
Orinoquia-Andina         -5.12334 -18.6574   8.41069  0.852164
Pacifica-Andina          -6.78912 -13.5628  -0.01545  0.049231
Pacifica-Caribe           1.56509  -5.9721   9.10230  0.980561
Orinoquia-Caribe          3.23087 -10.5362  17.00796  0.960027
Pacifica-Orinoquia       -1.66578 -15.9483  12.61671  0.995643
\end{lstlisting}

\vspace{0.3cm}
\textbf{Interpretación}:
\begin{itemize}
  \item Caribe vs Andina: p-adj $= 0.0018 < 0.05$ $\Rightarrow$ \alert{Diferencia significativa} (Caribe 8.4 puntos menor)
  \item Pacífica vs Andina: p-adj $= 0.049 < 0.05$ $\Rightarrow$ \alert{Diferencia significativa} (Pacífica 6.8 puntos menor)
  \item Otras comparaciones: No significativas
\end{itemize}
\end{frame}

\begin{frame}[fragile]{Visualización de Resultados Tukey}
\begin{lstlisting}
library(ggplot2)
library(broom)

# Extraer resultados de Tukey en formato tidy
tukey_df <- tidy(tukey_resultado)

# Grafico de forest plot
ggplot(tukey_df, aes(x = estimate, y = comparison)) +
  geom_point(size = 3) +
  geom_errorbarh(aes(xmin = conf.low, xmax = conf.high),
                 height = 0.2) +
  geom_vline(xintercept = 0, linetype = "dashed",
             color = "red", size = 1) +
  labs(x = "Diferencia de medias (IC 95% ajustado)",
       y = "Comparacion",
       title = "Tukey HSD: Comparaciones Multiples",
       subtitle = "Diferencias significativas si IC no incluye 0") +
  theme_minimal()
\end{lstlisting}
\end{frame}

% ============================================================================
\section{Curva de Potencia para ANOVA}
% ============================================================================

\begin{frame}{Tamaño del Efecto en ANOVA: $f$ de Cohen}
\begin{definition}[f de Cohen]
\[
f = \sqrt{\frac{\sum_{i=1}^k \pi_i (\mu_i - \mu)^2}{\sigma^2}}
\]
donde $\pi_i = n_i / N$ (proporción del grupo $i$), $\mu = \sum \pi_i \mu_i$ (media global).
\end{definition}

\vspace{0.3cm}
\textbf{Interpretación (Cohen, 1988)}:
\begin{itemize}
  \item $f \approx 0.10$: efecto \alert{pequeño}
  \item $f \approx 0.25$: efecto \alert{mediano}
  \item $f \approx 0.40$: efecto \alert{grande}
\end{itemize}

\vspace{0.3cm}
\textbf{Nota}: $f$ mide cuánto varían las medias de los grupos en relación a la variabilidad dentro de los grupos.
\end{frame}

\begin{frame}[fragile]{Calcular Potencia para ANOVA}
\begin{lstlisting}
library(pwr)

# Potencia para ANOVA de una via
# Parametros:
# k = numero de grupos
# n = tamano de muestra POR GRUPO
# f = f de Cohen
# sig.level = alpha

# Ejemplo: 5 grupos, n = 100 por grupo, f = 0.25, alpha = 0.05
potencia_anova <- pwr.anova.test(k = 5, n = 100, f = 0.25,
                                 sig.level = 0.05)

print(potencia_anova)
cat("\nPotencia:", round(potencia_anova$power, 3))

# Calcular n requerido para potencia 0.80
n_requerido <- pwr.anova.test(k = 5, f = 0.25, sig.level = 0.05,
                              power = 0.80)

cat("\nn requerido por grupo:", ceiling(n_requerido$n))
\end{lstlisting}
\end{frame}

\begin{frame}[fragile]{Explorar Diferentes Escenarios}
\begin{lstlisting}
# Tabla de n requerido para diferentes efectos
efectos <- c(0.10, 0.25, 0.40)
k_grupos <- c(3, 5, 10)

resultados <- expand.grid(f = efectos, k = k_grupos)
resultados$n_por_grupo <- sapply(1:nrow(resultados), function(i) {
  ceiling(pwr.anova.test(k = resultados$k[i],
                         f = resultados$f[i],
                         sig.level = 0.05,
                         power = 0.80)$n)
})

# Tabla pivotada
library(tidyr)
tabla_n <- resultados %>%
  pivot_wider(names_from = k, values_from = n_por_grupo,
              names_prefix = "k=")

print(tabla_n)
\end{lstlisting}

\vspace{0.3cm}
\small
\textbf{Observación}: Efectos pequeños requieren muestras MUCHO más grandes.
\end{frame}

% ============================================================================
\section{Pruebas No Paramétricas}
% ============================================================================

\begin{frame}{¿Cuándo Usar Pruebas No Paramétricas?}
\textbf{Situaciones}:
\begin{enumerate}
  \item \textbf{Datos ordinales} (ej. escala Likert: 1=malo, 2=regular, ..., 5=excelente)
  \item \textbf{Distribuciones muy sesgadas} o con outliers extremos
  \item \textbf{Muestras pequeñas} donde no se puede verificar normalidad
  \item \textbf{Violación de supuestos} de las pruebas paramétricas
\end{enumerate}

\vspace{0.3cm}
\textbf{Ventajas}:
\begin{itemize}
  \item Menos supuestos (no requieren normalidad)
  \item Robustas a outliers
  \item Válidas para datos ordinales
\end{itemize}

\vspace{0.3cm}
\textbf{Desventajas}:
\begin{itemize}
  \item \alert{Menos potentes} que las paramétricas si los supuestos se cumplen
  \item Interpretación menos directa (trabajan con rangos, no con medias)
\end{itemize}
\end{frame}

\begin{frame}{Pruebas No Paramétricas Comunes}
\begin{table}
\centering
\small
\begin{tabular}{lll}
\toprule
\textbf{Situación} & \textbf{Paramétrica} & \textbf{No paramétrica} \\
\midrule
1 muestra & Prueba $t$ & Wilcoxon signed-rank \\[0.2cm]
2 muestras independientes & Prueba $t$ & Mann-Whitney-Wilcoxon \\[0.2cm]
2 muestras pareadas & Prueba $t$ pareada & Wilcoxon signed-rank \\[0.2cm]
$k$ muestras independientes & ANOVA & Kruskal-Wallis \\[0.2cm]
$k$ muestras pareadas & ANOVA repetidas & Friedman \\
\bottomrule
\end{tabular}
\end{table}

\vspace{0.3cm}
\textbf{Principio}: Las pruebas no paramétricas trabajan con \alert{rangos} en lugar de valores originales.
\end{frame}

\begin{frame}{Mann-Whitney-Wilcoxon (U de Mann-Whitney)}
\textbf{Uso}: Alternativa a la prueba $t$ de dos muestras independientes.

\vspace{0.3cm}
\textbf{Hipótesis}:
\[
H_0: \text{Las dos distribuciones son idénticas}
\]
\[
H_1: \text{Las distribuciones difieren en localización}
\]

\vspace{0.3cm}
\textbf{Método}:
\begin{enumerate}
  \item Combinar todas las observaciones y asignar rangos
  \item Calcular la suma de rangos para cada grupo
  \item Comparar con la distribución nula
\end{enumerate}

\vspace{0.3cm}
\textbf{Interpretación}: Si rechazamos $H_0$, concluimos que un grupo tiende a tener valores \alert{mayores} (o menores) que el otro, pero NO necesariamente medias diferentes.
\end{frame}

\begin{frame}[fragile]{Ejemplo: Mann-Whitney en R}
\begin{lstlisting}
# Comparar ingles entre genero (si la distribucion no es normal)

# Dividir por genero
hombres <- icfes_ni %>%
  filter(ESTU_GENERO == "M") %>%
  pull(MOD_INGLES_PUNT)

mujeres <- icfes_ni %>%
  filter(ESTU_GENERO == "F") %>%
  pull(MOD_INGLES_PUNT)

# Prueba de Mann-Whitney
resultado_mw <- wilcox.test(hombres, mujeres,
                            alternative = "two.sided")

print(resultado_mw)

cat("\nEstadistico W:", resultado_mw$statistic)
cat("\np-valor:", resultado_mw$p.value)

# Decision
if (resultado_mw$p.value < 0.05) {
  cat("\nHay diferencia significativa entre generos")
} else {
  cat("\nNo hay diferencia significativa")
}
\end{lstlisting}
\end{frame}

\begin{frame}[fragile]{Comparar Paramétrica vs No Paramétrica}
\begin{lstlisting}
# Prueba t
resultado_t <- t.test(hombres, mujeres)
cat("Prueba t:")
cat("\n  t =", resultado_t$statistic)
cat("\n  p-valor =", resultado_t$p.value)

# Mann-Whitney
resultado_mw <- wilcox.test(hombres, mujeres)
cat("\n\nMann-Whitney:")
cat("\n  W =", resultado_mw$statistic)
cat("\n  p-valor =", resultado_mw$p.value)

# Comparar conclusiones
cat("\n\nAmbas pruebas coinciden?",
    (resultado_t$p.value < 0.05) == (resultado_mw$p.value < 0.05))
\end{lstlisting}

\vspace{0.3cm}
\small
\textbf{Nota}: Si los datos son aproximadamente normales, ambas pruebas deberían dar resultados similares. Si difieren mucho, investigar outliers o asimetría.
\end{frame}

\begin{frame}{Kruskal-Wallis: Alternativa No Paramétrica a ANOVA}
\textbf{Uso}: Comparar $k \geq 3$ grupos cuando los supuestos de ANOVA no se cumplen.

\vspace{0.3cm}
\textbf{Hipótesis}:
\[
H_0: \text{Las } k \text{ distribuciones son idénticas}
\]
\[
H_1: \text{Al menos una distribución difiere en localización}
\]

\vspace{0.3cm}
\textbf{Método}:
\begin{enumerate}
  \item Asignar rangos a todas las observaciones (ignorando grupos)
  \item Calcular la suma de rangos para cada grupo
  \item Estadístico $H$ (similar a $\chi^2$):
\end{enumerate}
\[
H = \frac{12}{N(N+1)} \sum_{i=1}^k \frac{R_i^2}{n_i} - 3(N+1)
\]
donde $R_i$ = suma de rangos del grupo $i$.

\vspace{0.3cm}
Si $H_0$ es verdadera, $H \sim \chi^2(k-1)$ aproximadamente.
\end{frame}

\begin{frame}[fragile]{Ejemplo: Kruskal-Wallis en R}
\begin{lstlisting}
# Comparar ingles entre macro-regiones (alternativa no parametrica)

resultado_kw <- kruskal.test(MOD_INGLES_PUNT ~ MACRO_REGION,
                             data = icfes_ni)

print(resultado_kw)

cat("\nEstadistico H (chi-cuadrado):", resultado_kw$statistic)
cat("\ngl:", resultado_kw$parameter)
cat("\np-valor:", resultado_kw$p.value)

# Decision
if (resultado_kw$p.value < 0.05) {
  cat("\nHay diferencia significativa entre regiones")
  cat("\n(al menos una difiere)")
} else {
  cat("\nNo hay diferencia significativa")
}
\end{lstlisting}
\end{frame}

\begin{frame}[fragile]{Comparaciones Post-hoc para Kruskal-Wallis}
\begin{lstlisting}
# Kruskal-Wallis no incluye comparaciones multiples automaticas
# Opcion 1: Usar Dunn test (paquete dunn.test)
library(dunn.test)

dunn_resultado <- dunn.test(icfes_ni$MOD_INGLES_PUNT,
                            icfes_ni$MACRO_REGION,
                            method = "bonferroni")
# method: "bonferroni", "holm", "hochberg", "bh", etc.

# Opcion 2: Pairwise Wilcoxon con ajuste
pairwise_resultado <- pairwise.wilcox.test(
  icfes_ni$MOD_INGLES_PUNT,
  icfes_ni$MACRO_REGION,
  p.adjust.method = "bonferroni"
)

print(pairwise_resultado)
\end{lstlisting}

\vspace{0.3cm}
\small
\textbf{Nota}: El ajuste de Bonferroni es conservador pero simple. Otras opciones: Holm, Hochberg, Benjamini-Hochberg (FDR).
\end{frame}

\begin{frame}{¿Cuándo Usar Paramétrica vs No Paramétrica?}
\textbf{Guía de decisión}:

\vspace{0.3cm}
\begin{enumerate}
  \item \textbf{Verificar supuestos}:
  \begin{itemize}
    \item QQ-plot, Shapiro-Wilk para normalidad
    \item Levene para homocedasticidad
    \item Tamaño de muestra
  \end{itemize}

  \item \textbf{Si los supuestos se cumplen}: Usar prueba paramétrica
  \begin{itemize}
    \item Mayor potencia
    \item Interpretación más directa (medias)
  \end{itemize}

  \item \textbf{Si los supuestos NO se cumplen}:
  \begin{itemize}
    \item Intentar transformación (log, raíz cuadrada)
    \item Si persiste, usar prueba no paramétrica
  \end{itemize}

  \item \textbf{Si los datos son ordinales}: Siempre usar no paramétrica
\end{enumerate}

\vspace{0.3cm}
\textbf{Consejo}: Con muestras grandes ($n \geq 30$ por grupo), la prueba $t$ y ANOVA son \alert{robustas} a desviaciones moderadas de la normalidad (TLC).
\end{frame}

% ============================================================================
\section{Resumen de la Unidad 3 (hasta ahora)}
% ============================================================================

\begin{frame}{Tabla Resumen: ¿Qué Prueba Usar?}
\begin{table}
\centering
\tiny
\begin{tabular}{p{2.5cm}p{2cm}p{2.5cm}p{2.5cm}}
\toprule
\textbf{Situación} & \textbf{Datos} & \textbf{Paramétrica} & \textbf{No paramétrica} \\
\midrule
1 muestra & Cuantitativos & Prueba $t$ (1 muestra) & Wilcoxon signed-rank \\[0.2cm]
2 muestras independientes & Cuantitativos & Prueba $t$ (2 muestras) & Mann-Whitney-Wilcoxon \\[0.2cm]
2 muestras pareadas & Cuantitativos & Prueba $t$ pareada & Wilcoxon signed-rank \\[0.2cm]
$k$ muestras independientes & Cuantitativos & ANOVA + Tukey HSD & Kruskal-Wallis + Dunn \\[0.2cm]
2 proporciones & Categóricos & Prueba $z$ para 2 prop. & (Próxima sesión: $\chi^2$) \\[0.2cm]
$k$ proporciones & Categóricos & (Próxima sesión: $\chi^2$) & (Próxima sesión: $\chi^2$) \\
\bottomrule
\end{tabular}
\end{table}

\vspace{0.3cm}
\small
\textbf{Próxima sesión (después del parcial)}: Pruebas $\chi^2$ para independencia y bondad de ajuste.
\end{frame}

\begin{frame}{Flujo de Decisión}
\begin{center}
\begin{tikzpicture}[
  node distance=1.2cm and 1.5cm,
  decision/.style={diamond, draw, fill=blue!20, text width=2.5cm, align=center, inner sep=1pt, minimum height=0.8cm},
  block/.style={rectangle, draw, fill=green!20, text width=3cm, align=center, rounded corners, minimum height=0.7cm},
  line/.style={draw, -{Latex[length=2mm]}}
]

  \node [block] (inicio) {?`Cu\'antos grupos?};

  \node [decision, below=1cm of inicio] (grupos) {1, 2, o $>$2?};

  \node [block, below left=1.2cm and 2.5cm of grupos] (un_grupo) {1 grupo:\\ Prueba $t$ (1 muestra)};

  \node [decision, below=1.8cm of grupos] (dos_grupos) {?`Independientes?};
  \node [block, below left=1cm and 1.5cm of dos_grupos] (t_dos) {Prueba $t$\\ (2 muestras)};
  \node [block, below right=1cm and 1.5cm of dos_grupos] (t_pareada) {Prueba $t$\\ pareada};

  \node [block, below right=1.2cm and 2.5cm of grupos] (anova) {$>$2 grupos:\\ ANOVA};

  \path [line] (inicio) -- (grupos);
  \path [line] (grupos) -- node[left, sloped] {\tiny 1} (un_grupo);
  \path [line] (grupos) -- node[above] {\tiny 2} (dos_grupos);
  \path [line] (grupos) -- node[right, sloped] {\tiny $>$2} (anova);
  \path [line] (dos_grupos) -- node[above] {\tiny S\'i} (t_dos);
  \path [line] (dos_grupos) -- node[above] {\tiny No} (t_pareada);

\end{tikzpicture}
\end{center}

\vspace{0.3cm}
\small
Si los supuestos no se cumplen, cambiar a la prueba no paramétrica correspondiente.
\end{frame}

% ============================================================================
\section{Conclusiones}
% ============================================================================

\begin{frame}{Resumen de la Sesión}
\textbf{Hoy aprendimos}:
\begin{enumerate}
  \item \textbf{Prueba $t$ pareada}: Para muestras dependientes
  \begin{itemize}
    \item Trabajar con las diferencias: $t = \frac{\bar{d}}{s_d/\sqrt{n}}$
  \end{itemize}

  \item \textbf{ANOVA de una vía}: Comparar $\geq 3$ grupos
  \begin{itemize}
    \item $F = \text{MSB}/\text{MSW}$
    \item Supuestos: normalidad, homocedasticidad, independencia
  \end{itemize}

  \item \textbf{Tukey HSD}: Comparaciones múltiples post-hoc
  \begin{itemize}
    \item Controla FWER al nivel $\alpha$
  \end{itemize}

  \item \textbf{Potencia en ANOVA}: $f$ de Cohen, tamaño de muestra

  \item \textbf{Pruebas no paramétricas}:
  \begin{itemize}
    \item Mann-Whitney (alternativa a $t$ de 2 muestras)
    \item Kruskal-Wallis (alternativa a ANOVA)
  \end{itemize}
\end{enumerate}

\vspace{0.3cm}
\textbf{Próxima sesión}: ¡Parcial! (cubre Sesiones 1-7)
\end{frame}

\begin{frame}{Preparación para el Parcial}
\textbf{Temas a repasar}:
\begin{enumerate}
  \item \textbf{Sesiones 1-3}: Estadística descriptiva, visualización, muestreo, correlación
  \item \textbf{Sesiones 4-5}: Estimación puntual, IC (medias, varianzas, proporciones), tamaño de muestra
  \item \textbf{Sesión 6}: Pruebas de hipótesis, p-valor, errores tipo I/II, potencia, d de Cohen
  \item \textbf{Sesión 7 (hoy)}: Prueba $t$ pareada, ANOVA, Tukey, no paramétricas
\end{enumerate}

\vspace{0.3cm}
\textbf{Habilidades clave}:
\begin{itemize}
  \item Interpretar output de R (\code{t.test()}, \code{aov()}, \code{summary()})
  \item Verificar supuestos (normalidad, homocedasticidad)
  \item Elegir la prueba adecuada según la situación
  \item Reportar resultados completos (estadístico, gl, p-valor, IC, tamaño del efecto)
\end{itemize}

\vspace{0.3cm}
\textbf{Recurso}: Repasar ejercicios de las sesiones anteriores con datos ICFES.
\end{frame}

\begin{frame}{Ejercicios para Practicar}
\textbf{Con los datos de ICFES}:
\begin{enumerate}
  \item Comparar \variable{MOD\_RAZONA\_CUANTITAT} antes/después de una intervención (simular datos pareados)
  \item Realizar ANOVA comparando inglés entre 5 departamentos (elegir 5 con $n$ grande)
  \item Aplicar Tukey HSD y reportar qué departamentos difieren
  \item Verificar supuestos de ANOVA (QQ-plot, Levene)
  \item Comparar inglés entre 3 niveles socioeconómicos usando Kruskal-Wallis
  \item Calcular potencia de ANOVA con $k=5$, $n=100$, $f=0.30$
  \item \textbf{Desafío}: Crear un reporte completo de análisis comparativo (descriptivos + visualización + prueba + interpretación)
\end{enumerate}

\vspace{0.3cm}
\textbf{Lectura}: Anderson et al., Capítulos 10-11 (Comparación de Medias, ANOVA)
\end{frame}

\begin{frame}
\begin{center}
\Large
¡Gracias por su atención!

\vspace{0.5cm}
\normalsize
Preguntas y discusión

\vspace{1cm}
\large
¡Mucho éxito en el parcial!
\end{center}
\end{frame}

\end{document}
